\section{Conclusions}

The EESM voltage surface is controlled by a simple dimensional fact: a linear
map from three current coordinates to two stator-voltage components has a null
space.  Its Gram matrix therefore has two positive eigenvalues and one zero
eigenvalue in the nondegenerate case.  The resulting surface is an oblique
elliptic cylinder assembled from congruent fixed-excitation ellipses, not an
ellipsoid.

Torque adds an indefinite bilinear geometry: nonzero levels are hyperbolic
cylinders and the zero level is a pair of planes.  Copper loss, by contrast,
has a positive-definite metric and forms a genuine ellipsoid.  This distinction
supports a minimum-loss tangency interpretation and a generalized eigenvalue
solution after loss whitening.

Finally, a voltage cylinder does not compactly bound torque.  Voltage-only MTPV
is generically unbounded in the linear model; a finite high-speed optimum
requires explicit thermal, stator-current, excitation-current, or saturation
constraints.  Numerical boxes must never be mistaken for those physical
constraints.  Rank, null spaces, and active-set checks remain useful even when
nonlinear data warp the ideal quadrics.
